\documentclass[12pt]{article}

\usepackage{amsfonts,amsmath,amssymb}
\DeclareMathAlphabet{\mathpzc}{OT1}{pzc}{m}{it}
\usepackage[italian]{babel}
\usepackage{caption}
\usepackage{enumitem}
\usepackage[a4paper,left=1.5cm,right=1.5cm,top=2.5cm,bottom=2.5cm]{geometry}
\usepackage{float}
\pagestyle{plain}
\usepackage{geometry}
\usepackage{graphicx}
\usepackage{setspace}
\usepackage{hyperref}
\usepackage{mathtools}
\usepackage{nicematrix}
\usepackage{subfigure}
\usepackage{titling}
\onehalfspacing 

\bibliographystyle{alpha}

\begin{document}
%%%%%%%%%%% TITLE %%%%%%%%%%%%%%%%%
\noindent LITAL1300 - preparazione all'esame orale\\[-2mm]
\rule{\linewidth}{0.5pt}

\begin{center}
    {\large LITAL1300 -- preparazione all'esame orale} \\
    \bigskip 
    Issambre L'Hermite Dumont \bigskip \\
\end{center}
\hrule
\bigskip
%%%%%%%%%%% END TITLE %%%%%%%%%%%%%

\section*{Video: Presentazione della Ferrari 812}
\url{https://www.youtube.com/watch?v=DggqGPGg6UU}
Ho deciso di guardare il video di presentazione della Ferrari 812 perché per me è una delle Ferrari moderne più riuscite insieme alla Ferrari 488 Pista. Inoltre, è legata alla mia altra attività perché discende dalla Ferrari 250, di cui la 250 GTO è la versione da corsa. Fa parte della linea delle Ferrari con motore anteriore e è per questo che l'ho scelta al posto della 488 Pista.\\ 
La Ferrari 812 è stata presentata nel 2017 e fa parte di una serie di Ferrari con motore V12 anteriore centrale. Il modello che l'ha preceduta era la Ferrari F12, mentre il modello successivo, lanciato nel 2024, si chiama Ferrari 12 Cilindri.\\
E una vettura que porta all'estremo il concetto del motore 12 cilindri perché sviluppa oltre 800 cavalli a 7900 giri al minuti. Anche questa vettura è dotata di un telaio ad alte prestazioni grazie all'introduzione delle quattro ruote sterzanti. La 812 è una vettura a trazione posteriore. Accelera da 0 a 100 km/h in meno di 3 secondi e raggiunge una velocità massima di 340 km/h.\\
La 812 si chiama "812" perché  è il rapporto tra cilindrata e potenza: $6 496/800 = 8,12 cm^3/CV$.\\
La 812 è stata progettata da Flavio Manzonni ed è stata concepita per garantire prestazioni elevate, con un'aerodinamica studiata per avere miglior draghi e quidi maggior velocità, pur mantenendo un aspetto estetico slanciato.\\ 
È stata progettata con numerosi elementi che guidano il flusso d'aria, come le prese d'aria anteriori. Ma anche nella parte posteriore, con l'estrattore d'aria in carbonio. Questa vettura presenta molti componenti in carbonio: il carbonio è leggero e resistente, è perché è utilizzato. Apprezzo molto tutti questi piccoli dettagli in carbonio.\\
Durante la mia ricerca nel primo quadrimestre, avevo detto che il rombo dei motori delle auto era come musica, e penso che la 812 sia una delle auto il cui motore suona come musica.\\
La Ferrari nel video è presentata in rosso Ferrari e, nonostante mi piaccia molto il rosso Ferrari per ciò che rappresenta, lo trovo un po’ deludente e spento su questa vettura. Preferirei un rosso più intenso e brillante o metallizzato, come il rosso Magma. \\
\section*{Articolo: La Ferrari più costosa del mondo}
\url{https://www.auto.it/news/attualita/2023/11/15-6850759/la_ferrari_piu_costosa_mai_venduta_all_asta_e_la_250_gto_del_1962}\\
Ho scelto questo articolo perché, secondo me, la Ferrari 250 GTO è la Ferrari più bella mai progettata e anche la più iconica. Di questa Ferrari, prodotta tra il 1962 e il 1964, ne esistono solo 39 esemplari: 33 "normali" con un motore da 3,0 litri, 3 con un motore da 4,0 litri e 3 "Typo 64" con un carrozzeria modificata. Per poterne acquistare una, bisognava essere in confidenza con lo stesso Enzo Ferrari. Il prezzo all’epoca era di $ 18.000\$ $, che oggi equivalgono a circa $ 200.000\$ $.\\
La 250 GTO era un'auto da corsa dalle prestazioni eccezionali per l'epoca. Ha vinto numerose corse ed è stata campionessa per tre anni consecutivi nella categoria Gran Turismo.\\
Nel 2023, una Ferrari 250 GTO del 1962 è stata venduta a New York per 51,7 milioni di dollari, anche se la sua stima era di 60 milioni. Si è trattato di un record per quell’anno. Nonostante il prezzo molto elevato, potrebbe ancora aumentare di valore. È un’auto speciale perché è l'unica del 1962 ad essere stata utilizzata in gara dalla Scuderia Ferrari. L'auto con il telaio n. 3765 vinse nella categoria Gran Turismo alla 1.000 chilometri del Nürburgring del 1962. Tuttavia, non riuscì a terminare la 24 Ore di Le Mans dello stesso anno a causa di problemi di surriscaldamento.\\
La vettura era equipaggiata con un motore V12 da 4,0 litri. Ma dopo Le Mans, il motore è stato sostituito con il motore base da 3,0 litri. Nel 1985 è arrivata negli Stati Uniti, dove ha partecipato a corse di auto storiche. Nel 2011 ha vinto un Concorso d'Eleganza.\\
\section*{10 parole}
\begin{enumerate}
    \item \textbf{Telaio}: È lo scheletro dell'auto, e quello che tiene in posizione le ruote, il motore e tutto il resto.
    \item \textbf{Draghi}: La resistenza del aria
    \item \textbf{Prestazioni}: Quando un'auto va veloce e funziona bene, è una questione di prestazioni
    \item \textbf{Surriscaldamento}: Quando il motore è troppo caldo
    \item \textbf{Sterzanti}: Que fare girare l'auto. Le \textbf{quattro ruote sterzanti} significa che anche le ruote posteriori fare girare l'auto.
    \item \textbf{Pittura}: Il colore dell'auto
    \item \textbf{Contagiri}: È uno strumento che mostra a quanti giri al minuto sta girando il motore
    \item \textbf{Altisonante}: Molto grande, impressionante
    \item \textbf{Fruitibilita}: La facilità con cui si può usare la vettura
    \item \textbf{Controsterzo}: Quando si fa girare le ruote in una direzione per far girare l'auto in un'altra direzione, si chiama controsterzo.
\end{enumerate}
\end{document}